\section{Maximal consistent sets}
\label{sec:prop-metatheory-02}

\begin{definition}[Consistency]
A set \(\Gamma\subseteq\Fm_{\CPL}\) is \emph{consistent} iff
\(\Gamma\notderives_{\ND_{\CPL}}\bottom\).
\end{definition}

\begin{definition}[Maximal consistent set]
A \emph{maximal consistent set} is a consistent set
\(\Delta\subseteq\Fm_{\CPL}\) such that
\(\Delta\cup\set{\varphi}\) is inconsistent for every
\(\varphi\notin\Delta\).
\end{definition}

\begin{lemma}[Closure properties]
\label{lem:prop-mcs}
If \(\Delta\) is maximal consistent, then:
\begin{enumerate}
  \item if \(\Delta\derives_{\ND_{\CPL}}\varphi\), then
        \(\varphi\in\Delta\);
  \item exactly one of \(\varphi\) and \(\neg\varphi\) belongs to
        \(\Delta\);
  \item \(\varphi\land\psi\in\Delta\) iff
        \(\varphi\in\Delta\) and \(\psi\in\Delta\);
  \item \(\varphi\lor\psi\in\Delta\) iff
        \(\varphi\in\Delta\) or \(\psi\in\Delta\); and
  \item \(\varphi\to\psi\in\Delta\) iff
        \(\varphi\notin\Delta\) or \(\psi\in\Delta\).
\end{enumerate}
\end{lemma}
\begin{proof}
For 1, suppose \(\Delta\derives\varphi\) but \(\varphi\notin\Delta\).
Maximality makes \(\Delta\cup\set{\varphi}\) inconsistent.  Substitute the
derivation of \(\varphi\) for every open use of the added assumption by
\cref{lem:nd-substitution}; this yields a derivation of \(\bottom\) from
\(\Delta\), contrary to consistency.

The two formulas in 2 cannot both belong to \(\Delta\), by \(\neg E\).  If
\(\varphi\notin\Delta\), maximality yields a derivation \(D\) of
\(\bottom\) whose open assumptions lie in \(\Delta\cup\set{\varphi}\).
Because \(\Delta\) is consistent, at least one open occurrence of
\(\varphi\) remains in \(D\).  Mark all open occurrences of \(\varphi\) in
\(D\) with a common fresh label \(i\), by the merging relabeling of the
derivation definition, and apply \(\neg I\) with label \(i\) at the root,
discharging every such occurrence, to obtain
\(\Delta\derives\neg\varphi\); then 1 gives \(\neg\varphi\in\Delta\).
This proves 2.

For 3, conjunction elimination proves the forward implication and conjunction
introduction proves the reverse implication, using 1.  For 4, disjunction
introduction proves the reverse implication.  If
\(\varphi\lor\psi\in\Delta\) while neither disjunct belongs to \(\Delta\),
then 2 gives \(\neg\varphi,\neg\psi\in\Delta\).  Disjunction elimination,
using each negation to derive \(\bottom\) in its corresponding branch, derives
\(\bottom\) from \(\Delta\), a contradiction.

For 5, if \(\varphi\to\psi\in\Delta\) and \(\varphi\in\Delta\), implication
elimination and 1 give \(\psi\in\Delta\).  Conversely, if
\(\psi\in\Delta\), assume \(\varphi\) with label \(i\) and take \(\psi\)
as an open assumption; \(\land I\) gives \(\varphi\land\psi\), and
\(\land E_2\) gives \(\psi\); then \(\to I\) discharges the single open
occurrence marked \(i\), concluding \(\varphi\to\psi\) with open
assumptions in \(\Delta\).  The discharged occurrence is genuinely
present, and only displayed rules occur.  If \(\varphi\notin\Delta\), then
\(\neg\varphi\in\Delta\) by 2; from an assumed \(\varphi\), derive
\(\bottom\), then \(\psi\), and discharge the assumption, which occurs
open in that derivation.  In either case 1
places \(\varphi\to\psi\) in \(\Delta\).
\end{proof}

\begin{lemma}[Lindenbaum lemma]
\label{lem:prop-lindenbaum}
Every consistent subset of \(\Fm_{\CPL}\) is contained in a maximal
consistent set.
\end{lemma}
\begin{proof}
The set \(\Fm_{\CPL}\) is countable; enumerate it as
\(\varphi_0,\varphi_1,\ldots\).  Starting with a consistent set
\(\Gamma_0=\Gamma\), define
\[
 \Gamma_{n+1}=
 \begin{cases}
   \Gamma_n\cup\set{\varphi_n},&
     \text{if this set is consistent},\\
   \Gamma_n,&\text{otherwise}.
 \end{cases}
\]
Let \(\Delta=\bigcup_{n\in\N}\Gamma_n\).  If \(\Delta\) were inconsistent,
a finite derivation of \(\bottom\) would use only finitely many assumptions,
all lying in some \(\Gamma_n\), contradicting its consistency.  Thus
\(\Delta\) is consistent.  If \(\varphi_n\notin\Delta\), it was rejected at
stage \(n+1\), so \(\Gamma_n\cup\set{\varphi_n}\), and hence
\(\Delta\cup\set{\varphi_n}\), is inconsistent.  Therefore \(\Delta\) is
maximal consistent.
\end{proof}

\begin{lemma}[Truth lemma]
\label{lem:prop-truth}
Let \(\Delta\) be maximal consistent and define a valuation \(v_\Delta\) by
\[
  v_\Delta(p_n)=\True\quad\text{iff}\quad p_n\in\Delta.
\]
Then, for every propositional formula \(\varphi\),
\[
  v_\Delta\entails\varphi
  \quad\text{iff}\quad
  \varphi\in\Delta.
\]
\end{lemma}
\begin{proof}
Proceed by structural induction.  The variable case is the definition of
\(v_\Delta\), and \(\bottom\notin\Delta\) by consistency.  The negation case
uses \cref{lem:prop-mcs}(2) and the truth clause for \(\neg\).  The
conjunction, disjunction, and implication cases use respectively clauses 3,
4, and 5 of that lemma together with the corresponding recursive truth
clauses.  These are all constructors in the grammar, so the induction is
complete.
\end{proof}

